\begin{pastexam}[label={px:ncku-112-prog-10}]{成大資訊聯招 112 學年度 第 10 題}
（複選）

{[5\%]} Please identify which of the following statement(s) is (are) true.
\begin{enumerate}[label=(\Alph*)]
\item Tries outperform digital search trees in terms of the number of comparison operations performed.
\item Compared with tries, compressed tries take less storage overheads.
\item Patricia outperforms compressed tries in terms of storage overheads.
\item Both Patricia and digital search trees take the identical complexity of storage overheads.
\item Patricia outperforms digital search trees in terms of the number of comparison operations
performed.
\end{enumerate}
\end{pastexam}

\begin{pastsoln}
答案：(A)(B)(C)(D)(E)，全對。

「comparison operations」指的是整個鍵的比較（一次比較要比完 $b$ 個位元，是貴的操作）；看單一位元決定往左
往右不算。 資結書第~11.5.4~節的比較表（$n$ 個 $b$ 位元的鍵）：
\begin{center}
\small
\begin{tabular}{lccc}
\toprule
& 節點數 & 搜尋時比幾次整個鍵 & 搜尋時看幾個位元 \\
\midrule
數位搜尋樹 & $n$ & 最多 $b + 1$ & 最多 $b$ \\
二元 trie & 分支節點最多 $nb$，加 $n$ 個元素節點 & $1$ & 最多 $b$ \\
壓縮二元 trie & $2n - 1$ & $1$ & 最多 $\min(b, n - 1)$ \\
Patricia & $n$（含標頭） & $1$ & 最多 $\min(b, n - 1)$ \\
\bottomrule
\end{tabular}
\end{center}

\begin{itemize}
\item (A) 對。 數位搜尋樹每經過一個節點都要把整個鍵和 $k$ 比一次，最多 $b + 1$ 次；trie 一路只看位元，走到
元素節點才比\textbf{一次}整個鍵。
\item (B) 對。 壓縮 trie 把只有一個子節點的分支節點拿掉，節點數從最多約 $nb$ 降到 $2n - 1$，儲存的額外負擔
較少。
\item (C) 對。 Patricia 再把元素節點也拿掉、鍵放進分支節點，只要 $n$ 個節點，比壓縮二元 trie 的 $2n - 1$ 少。
\item (D) 對。 兩者都是 $n$ 個節點，每個節點一個鍵、兩個指標（Patricia 多一個 $\mathit{bit}$ 欄位），空間都是
$\Theta(n)$ 個節點。
\item (E) 對。 Patricia 和 trie 一樣只在最後比一次整個鍵，數位搜尋樹最多比 $b + 1$ 次。
\end{itemize}
\end{pastsoln}
